\section{导读：为什么机器人需要基座模型}

本节先建立整期的问题。陈建宇的核心判断是，过去机器人常常是“100 个场景、100 个任务、100 套专用模型”，这种方式很难 scale。大语言模型和多模态模型的成功，让机器人领域看到另一条路：能不能为机器人训练一个更通用的 foundation model，让它跨任务、跨本体、跨场景泛化。

\begin{importantbox}{本期核心命题}
VLA，即 Vision-Language-Action Model，不只是给机器人接一个语言模型，而是让视觉、语言和动作进入同一个可扩展的端到端系统。真正的机器人基座模型，要从“利用现有基础模型做机器人”走向“为机器人预训练基础模型”。
\end{importantbox}

\begin{figure}[H]
\centering
\includegraphics[width=0.96\textwidth]{ch03-00693.jpg}
\caption{How robotic learning works now? 传统机器人学习通常围绕任务、数据、模型和执行闭环展开。投屏帧，约 00:11:33。}
\end{figure}

\begin{knowledgebox}{读图：从专用模型到通用模型的压力}
图中展示的不是单篇论文，而是整期课的背景：机器人学习长期依赖特定任务、特定数据和特定部署。陈建宇强调，这种方式面对家庭、工厂和人形机器人时不够通用，因而需要基座模型路线。
\end{knowledgebox}

\subsection{机器人 AI 的两阶段}

本节把路线分清楚。第一阶段是 leveraging foundation models in robotics，即把 LLM、VLM、Code LM 等现有基础模型拿来替代机器人系统中的某个模块。第二阶段是 pretraining foundation models for robotics，即直接为机器人训练能够输出动作的模型。前者是拼装和借力，后者才是真正的机器人基座模型。

\begin{figure}[H]
\centering
\includegraphics[width=0.90\textwidth]{ch05-01297.jpg}
\caption{Foundation models for robotics：从 LLM、VLA 到 Actuation 的模块关系。投屏帧，约 00:21:37。}
\end{figure}

\begin{knowledgebox}{读图：机器人三板块与基础模型替代}
图中把 Perception、Actuation 和 LLM 放在一起。传统机器人通常拆成规划、感知和执行；大模型浪潮先替代 Planning，再替代 Perception，最后才试图让 Action 也进入统一模型。
\end{knowledgebox}

\begin{knowledgebox}{术语消化：VLA、VLM、LLM}
\noindent\begin{tabularx}{\textwidth}{p{0.18\textwidth}p{0.34\textwidth}X}
\toprule
术语 & 含义 & 本课中的作用 \\
\midrule
LLM & Large Language Model，大语言模型 & 先替代或增强机器人规划。 \\
VLM & Vision-Language Model，视觉语言模型 & 用于感知、场景理解和视觉反馈。 \\
VLA & Vision-Language-Action Model，视觉语言动作模型 & 直接把视觉、语言和动作连接到端到端控制。 \\
Robot Foundation Model & 机器人基座模型 & 目标是跨任务、跨本体和跨场景泛化。 \\
\bottomrule
\end{tabularx}
\end{knowledgebox}

\subsection{本章小结}

VLA 的出现，是机器人从模块拼装走向通用基座模型的关键线索。整期课的主线，是从 LLM/VLM 借力，到 RT-2/OpenVLA/HiRT/pi0/GR00T N1，再到 diffusion policy、RDT 和未来统一理解预测模型。

